\documentclass[11pt]{article}
\usepackage[margin=1in]{geometry}\usepackage{booktabs,hyperref}
\title{Reproduction of ``A Stochastic Rounding-Enabled Low-Precision Floating-Point MAC for DNN Training'' (DATE 2024)}
\author{Samagna Potham}\date{\today}
\begin{document}\maketitle
\section{What was implemented}
FP8 (E5M2) exact multiplier, FP12 (E6M5) accumulator with round-to-nearest, lazy stochastic rounding and the paper's eager stochastic
rounding adder (Sticky Round + Round Correction), a Galois LFSR, and a complete MAC in Verilog. A bit-exact Python golden model,
a PyTorch emulator and a fused Numba CUDA kernel run the same arithmetic for training.
\section{Verification}
Adder RTL vs.\ golden model: 180{,}993 vectors at $r=18$ and 30{,}000 each at $r\in\{9,12,14,16\}$, zero mismatches (RN, lazy, eager);
MAC with LFSR: 4000 steps, zero mismatches. Paper's brute-force test (10000 pairs $\times$ 1000 random words): all results are one of the two
neighbouring FP12 values and no pair's round-up frequency is inconsistent with the ideal probability (exact binomial test, $p<10^{-6}$).
\section{Hardware cost (open-source proxy)}
Yosys + ABC mapped to the Nangate 45\,nm library, adders only (absolute numbers differ from the paper's Synopsys/28\,nm flow; compare ratios).
\begin{table}[h]\centering\begin{tabular}{lrrr}\toprule Design & $r$ & Area ($\mu$m$^2$) & Delay (ps)\\\midrule
RN E6M5 & -- & 895 & 1205\\ SR lazy E6M5 & 9 & 971 & 1281\\ SR eager E6M5 & 9 & 948 & 1188\\
SR lazy E6M5 & 18 & 1182 & 1431\\ SR eager E6M5 & 18 & 1145 & 1270\\\bottomrule\end{tabular}\end{table}
Eager is smaller and faster than lazy at both widths (area $-2$ to $-3\%$, delay $-7$ to $-11\%$); the paper reports larger gains.
Synthesis noise for identical logic is about $\pm7\%$ area. Energy, FP16/FP32 baselines and subnormal support were not evaluated.
\section{Training accuracy}
\begin{table}[h]\centering\caption{resnet20: paper vs.\ our emulation}\begin{tabular}{lrrrr}\hline
Configuration & Paper \% & Ours \% & Epochs & $\Delta$ \\\hline
FP32 baseline & 91.47 & -- & -- & -- \\
RN, FP16 acc. (E5M10) & 91.10 & -- & -- & -- \\
RN, BF16 acc. (E8M7) & 88.79 & -- & -- & -- \\
RN, E6M5 acc. & 83.03 & -- & -- & -- \\
SR eager, E6M5, r=9 & 43.11 & -- & -- & -- \\
SR eager, E6M5, r=12 & 89.34 & -- & -- & -- \\
SR eager, E6M5, r=16, no sub & 90.67 & -- & -- & -- \\
SR eager, E6M5, r=18, no sub & 91.39 & -- & -- & -- \\
\hline\end{tabular}\end{table}
\begin{table}[h]\centering\caption{vgg16: paper vs.\ our emulation}\begin{tabular}{lrrrr}\hline
Configuration & Paper \% & Ours \% & Epochs & $\Delta$ \\\hline
FP32 baseline & 93.46 & -- & -- & -- \\
RN, FP16 acc. & 93.06 & -- & -- & -- \\
SR eager, E6M5, r=18, no sub & 93.11 & -- & -- & -- \\
\hline\end{tabular}\end{table}
\begin{table}[h]\centering\caption{resnet50: paper vs.\ our emulation}\begin{tabular}{lrrrr}\hline
Configuration & Paper \% & Ours \% & Epochs & $\Delta$ \\\hline
FP32 baseline & 80.94 & -- & -- & -- \\
RN, FP16 acc. & 80.30 & -- & -- & -- \\
SR eager, E6M5, r=18, no sub & 80.33 & -- & -- & -- \\
\hline\end{tabular}\end{table}
\section{Deviations and limitations}
See \texttt{CHECKLIST.md}: assumed LFSR parameters, FP8 handling of subnormals/overflow, extra normalization path for far-path subtraction,
software FP16/BF16 baselines, no subnormal-support configurations, ResNet-50 preprocessing assumptions.
\end{document}
